\section{What was removed, narrowed, or retracted --- and why}
\label{sec:removed}

Every claim this project has withdrawn or restricted is listed here with the evidence that forced
the change. Nothing is dropped silently. A report that shows only its surviving results has
removed exactly the information a reader needs in order to calibrate the rest.

\subsection{Retracted: claims the data does not support}

\begin{description}[leftmargin=1.4em,style=nextline]
\item[{T1-kernel monotone ordering ($p = 1/120$).}]
  The perfectly ordered five-point decline was an artifact of \texttt{KA\_EXTRACT=strict}, a
  candidate filter that discards 66\% of 0.5B generations and 5\% of 14B generations --- a
  severity that correlates with the independent variable. Under the permissive policy Spearman
  $\rho$ moves from $-1.000$ to $-0.900$ and 3B breaks the order. \textbf{The endpoint contrast
  survives both policies and is retained}; the ordering claim and its $p$-value are withdrawn.

\item[{T1-kernel ``14B verifies nothing''.}]
  Recorded twice from a cell that was still running (0 verified at 108 and at 252 attempts). The
  finished cell has 2 verified in 338. 14B is near-zero, not zero, and the falsifier language now
  says so.

\item[{T3 ``frontier self-modification is one-shot''.}]
  The harness requests at most 12 strategies; 5 of 5 interpretable runs sit at 11--12 from round
  0 and 4 of 5 never grow. The observed plateau is the cap, not the model. Gains against a frozen
  procedure remain real and are reported.

\item[{T3 cap dose-response ($+0.482$ at unbounded).}]
  Came from runs carrying the EDQUOT signature ($Q_0 = 0.000$ and four to five all-zero rounds), a
  filesystem failure in which the grader raises and every candidate scores zero. Clean runs show
  no dose-response: $-0.046 / +0.021 / -0.009$.

\item[{``H100's \texttt{torch.compile} closed the headroom''.}]
  Offered as the explanation for a spike of scores at exactly $0.50$. It is false: the default
  scorer never touched the compiled baseline at the time. The spike is a plain-torch rewrite
  landing at parity, which is a different and more consequential fact.

\item[{Teacher injection improves compounding.}]
  First reported as $+0.100$ against $+0.051$ for the control. It does not survive dropping
  rounds whose lineage was severed by a resume defect, nor matching on round index, since the
  arms had unequal valid depth. Matched: control $+0.057\,[+0.020,+0.094]$ versus inject
  $+0.075\,[+0.035,+0.114]$; and on the completed pass-rate board, $+0.440$ versus $+0.445$.
  \textbf{Not detected.}

\item[{``Retrieval is the worst baseline''.}]
  Stated from one cell at $0.154$. The completed four-cell sweep puts retrieval at $0.154$,
  $0.329$, $0.385$ and $0.426$ --- a spread of $0.272$, against within-cell gaps between any two
  methods of at most $0.27$ and typically half that. It is the worst method in one cell and
  within $0.013$ of the best in another. Retrieval is high variance, not reliably weak, and the
  original claim was an artifact of reading the low end of that range as the value.

\item[{``76\% of all scores were exactly 0.50''.}]
  Not reproducible from any artifact. The defensible figure is that 75\% of \emph{non-zero}
  scores fall within $\pm 0.01$ of parity on the 29-task calibration set.
\end{description}

\subsection{Narrowed: claims that hold with a smaller scope}

\begin{description}[leftmargin=1.4em,style=nextline]
\item[{The compounding null ($n=57$ trajectories, \bfz{}$\,{=}\,5.6$).}]
  Valid, and narrowed twice. It is a null about acquiring \emph{correctness} rather than
  optimisation, and specifically about \textbf{best-of-$k$} correctness, which by construction
  cannot see a reliability gain. The pass-rate board finds a large, family-transferring effect on
  that axis. These are not contradictory: the null is valid and narrow, and its narrowness was
  invisible until a metric existed that could see past it.

\item[{T1-kernel shape across scale.}]
  Reported as a decline with 14B as the endpoint. The completed 32B sweep recovers to $2.99\%$
  from 14B's $0.00\%$ (Fisher $p = 0.015$), so the curve is a \textbf{U} with 14B as its
  \emph{minimum}. The 0.5B-versus-14B contrast is a contrast with the floor, not with the largest
  model.

\item[{Compile-gain ``$1.34\times$ median''.}]
  Correct and misleading alone. The distribution is bimodal: L1 $15.51\times$, L2 $1.26\times$,
  L3 $4.53\times$. The median is simply L2's value because L2 is the majority tier. Reported per
  tier throughout, because the argument that a plain-torch rewrite scores $0.50$ ``without
  compile being strong'' holds on L2 and fails on L1.
\end{description}

\subsection{Excluded: data that cannot answer the question it was collected for}

\begin{description}[leftmargin=1.4em,style=nextline]
\item[{Severed trajectories (Amendment 5).}]
  \texttt{lab\_compounding} and \texttt{lab\_weight\_rsi} restored round history on resume but not
  the trained adapter. Every resumed chunk continued the round numbering while restarting each
  arm from base weights. The signature is unambiguous: \texttt{trainC} collapses at exactly the
  round named by \texttt{resumed\_at}, in all 18 resumed cells with no exceptions. All affected
  rounds are excluded, and the generator refuses to report a cell whose artifact cannot evidence
  that its weights survived a resume.

\item[{Starved trajectories (Amendment 6).}]
  A trajectory whose self arm received fewer than two verified examples per round on average is
  reported as starved rather than as a measurement of the contrast. At the registered
  configuration every cell qualifies, so \textbf{the registered primary currently reports no
  number at all}, which is the correct outcome rather than a missing one.

\item[{The 32B probe's first reading.}]
  Its first 12 tasks are L1-heavy and gave $6.94\%$. Restricting every scale to those same 12
  tasks is the valid comparison; the full 29-task sweep gives $2.99\%$. The partial figure is
  recorded and not reported.
\end{description}

\subsection{Provisional: retained but not load-bearing}

\begin{description}[leftmargin=1.4em,style=nextline]
\item[{Probe-as-intervention ($+0.124\,[+0.029,+0.218]$).}]
  $n=12$ checkpoints, 5--12 tasks per run, against the weakest available comparator (uniform
  random). Reported in the appendix, not near the abstract.

\item[{T4 closed$\to$open.}]
  $n = 1$ per cell against a registered minimum of three seeds, with no interval. Exploratory.

\item[{Any A100 speed-scored board.}]
  Measured under a harness that silently rejected every non-PyTorch submission. Retained with the
  defect stated, never pooled with H100 sets.
\end{description}

\subsection{Known confounds carried into the results}

\begin{itemize}[leftmargin=1.4em]
\item \textbf{Cross-rung token budgets differ.} T1 and T2 generate at 2048 tokens; T3 and T5 at
  1200. Both resume from per-round checkpoints, so changing them mid-flight would mix budgets
  inside a single cell, which is worse than a consistent difference. Any cross-rung comparison
  carries a $1.7\times$ gap and says so.
\item \textbf{The held set is 5 tasks, not the 20 requested.} A slice does not raise when it runs
  past the end, so a 25-task pool with 20 taken for training returns whatever remains. At $k=10$
  this caps a round at 50 verified candidates, which bounds every per-round score regardless of
  its functional form.
\item \textbf{$n = 2$ trajectories per arm} on the pass-rate board. No interval is quoted on a
  two-point mean; the per-cell values are given instead.
\end{itemize}
